% Generated draft: requires mathematical and source review.
\documentclass[11pt]{article}
\usepackage[margin=25mm]{geometry}
\usepackage{fontspec}
\setmainfont{Latin Modern Roman}
\usepackage{amsmath,amssymb,mathtools,unicode-math}
\setmathfont{Latin Modern Math}
\usepackage{enumitem,needspace,xurl,hyperref}
\setlength{\parindent}{0pt}
\setlength{\parskip}{5pt}
\setlength{\emergencystretch}{3em}
\newcommand{\N}{\mathbb{N}}
\newcommand{\Z}{\mathbb{Z}}
\newcommand{\Q}{\mathbb{Q}}
\newcommand{\R}{\mathbb{R}}
\newcommand{\C}{\mathbb{C}}
\newcommand{\F}{\mathbb{F}}
\newcommand{\A}{\mathbb{A}}
\newcommand{\PP}{\mathbb{P}}
\newcommand{\E}{\mathbb{E}}
\newcommand{\eps}{\varepsilon}
\newcommand{\dd}{\,\mathrm{d}}
\begin{document}
\section*{UnsolvedMath ID: 30001776}
\textbf{Source problem code:} catalogue entry
\subsection{Definitions and mathematical statement}
\textbf{Source formulation.} Characterize the energy release rate for the very general crack geometries occurring in the Francfort--Marigo fracture model.

\textbf{Definitions and unresolved scope.} The Francfort--Marigo fracture model and its admissible crack geometries are named, but not mathematically defined, in the supplied record. In particular, the record specifies no ambient domain or dimension, elastic energy functional, displacement class, loading or boundary conditions, or admissible class of cracks. It also specifies neither a family of crack extensions nor a parameter or normalization with respect to which an energy release rate is to be defined. Thus the intended notion of energy release rate, including any derivative, one-sided limit, or alternative generalized notion, remains unspecified.

The requested characterization concerns the general crack geometries of the named model, rather than merely special geometries. However, the supplied record does not determine the hypotheses, quantifiers, or mathematical form of such a characterization. This is therefore a source-faithful research objective, not a fully specified proposition. A precise quantified formulation requires the missing model definitions and the intended meaning of energy release rate; no particular choices are imposed here.
\subsection{Short English statement}
The problem asks for a characterization of energy release rate for the general crack geometries allowed by the Francfort--Marigo fracture model. The supplied record does not define the model, admissible cracks, or rate, so a precise theorem statement cannot yet be extracted. Its open status is supplied curation, not an independently verified conclusion.
\subsection{Sources}
\begin{itemize}
\item Local registry: \texttt{data/unsolvedmath/problems.json}, with \texttt{catalog/problems.json} designated for additional existing records. Supplied complete source record for permanent UnsolvedMath ID \texttt{30001776}, source problem code \texttt{catalogue entry}, title \emph{Energy Release for General Crack Geometries}.
\item Supplied report citation: \emph{Mini-Workshop: Mathematical Models, Analysis, and Numerical Methods for Dynamic Fracture} (2011), Oberwolfach Reports, pp.~1207--1240, DOI \texttt{10.4171/owr/2011/22}. \url{https://doi.org/10.4171/owr/2011/22}. The record attributes the problem to the workshop-report collection, gives proposed year 2011, and names no individual proposer. Francfort and Marigo are identified through the model name, not as supplied problem proposers.
\item Supplied literature source labeled ``original OWR report'': \url{https://doi.org/10.4171/owr/2011/22}.
\item Supplied literature source labeled ``special crack geometries'': \url{https://doi.org/10.1007/s00332-010-9061-2}. No title, authors, or precise result for this reference are supplied.
\item Supplied literature assessment, checked \textbf{2026-08-21}: the general energy-release-rate characterization for Francfort--Marigo crack geometries was explicitly open in the report, and later work treats only special geometries. This is supplied dated curation derived from the OWR collection, not an independently verified literature assessment. The record labels the problem ``open'' and its statement status \texttt{corrected\_verified}; it supplies the verification claim that the official report confirms the requested characterization for general crack geometries. Neither that verification claim nor the time-sensitive status has been independently checked here.
\end{itemize}
\end{document}
